\section{Exact scalarization of two Cauchy pairs}
\label{sec:scalarization}

This section replaces a matrix conductivity by a scalar one while
preserving two specified Cauchy pairs exactly. This is a finite-field
statement: it preserves the full flux response to each of two prescribed
voltage inputs, not the matrix conductivity's full
Dirichlet-to-Neumann operator. Classical isotropic approximation in
$G$-convergence provides useful historical context
\cite{MarinoSpagnolo}, but no homogenization result is used below.
The exact boundary data follow from localized perturbations and a
strong limit selected by the Baire theorem.
The continuity-point method follows the Baire-category framework for
differential inclusions developed by Kirchheim~\cite{Kirchheim} and
used for elliptic equations by Astala, Faraco, and
Sz\'ekelyhidi~\cite[Lemmas~3.7--3.8]{AstalaFaracoSzekelyhidi}.
We give the complete
localized construction and limiting argument for the simultaneous
Cauchy-pair constraints needed here.

\begin{theorem}[Exact finite-field scalarization]
\label{thm:scalarization}
Let $U\subset\R^3$ be a bounded Lipschitz domain. Let $A$ be a
measurable symmetric matrix field satisfying
\[
 c_0\Id\leq A\leq C_0\Id\qquad\text{almost everywhere in }U,
 \qquad 0<c_0\leq C_0<\infty.
\]
Suppose the real-valued potentials $u_1,u_2\in H^1(U)$ satisfy
$\Div(A\nabla u_j)=0$. Assume that there is an open cover of a
full-measure subset of $U$ on whose members $A,u_1,u_2$ are smooth
and have the following local rank property. On each member, a fixed
subset of the potentials has everywhere independent gradients, and
every remaining potential is a constant affine combination of that
subset. The subset may have zero, one, or two elements; a zero-element
subset means that both potentials are constant there.

There exist constants $0<a<b<\infty$, a measurable scalar
$a\leq\gamma\leq b$, and $v_1,v_2\in H^1(U)$ such that
\begin{align}
 \Div(\gamma\nabla v_j)&=0,
 &v_j-u_j&\in H^1_0(U), \label{sc:conclusion-equations}\\
 \int_U\gamma\nabla v_j\cdot\nabla\psi\,\dd x
 &=\int_U A\nabla u_j\cdot\nabla\psi\,\dd x
 &&\text{for every }\psi\in H^1(U). \label{sc:conclusion-flux}
\end{align}
The bounds $a,b$ may be chosen using only $c_0,C_0$.
\end{theorem}

We use the Euclidean norm for vectors and the Frobenius norm for
matrices. Write $\mathrm{Sym}_3$ for the vector space of real symmetric
$3\times3$ matrices, with this norm. Gradient and flux matrices have
their fields as columns.
In particular, $E^tF$ is the matrix of their pointwise pairings.
The proof will not require a positive lower bound on the independent
gradients throughout $U$: such bounds are used only on individual
compactly contained coordinate boxes.

\subsection{An open class of finite laminates}

For positive triples $\alpha,\beta\in(0,\infty)^3$, a coordinate
direction $i$, and $0\leq\theta\leq1$, define their coordinate join
$J_i(\theta;\alpha,\beta)$ by
\begin{equation}
 \label{sc:join}
 [J_i]_i=\left(\frac{\theta}{\alpha_i}
                  +\frac{1-\theta}{\beta_i}\right)^{-1},
 \qquad
 [J_i]_\ell=\theta\alpha_\ell+(1-\theta)\beta_\ell
 \quad(\ell\ne i).
\end{equation}
Fix $0<a<b$. Let $\mathcal G^{\mathrm{diag}}$ consist of the
triples obtainable by a finite tree of these joins, starting from
scalar triples $(s,s,s)$ with $a<s<b$. A tree is evaluated in one
fixed coordinate frame. Let $\mathcal G$ be its spectral extension:
\[
 \mathcal G=
 \{Q\diag(\alpha)Q^t:
       Q\text{ orthogonal},\ \alpha\in\mathcal G^{\mathrm{diag}}\}.
\]
In particular, every matrix in $\mathcal G$ lies strictly between
$a\Id$ and $b\Id$.

\begin{lemma}[Openness and coverage]
\label{sc:open-class}
The set $\mathcal G^{\mathrm{diag}}$ is open in $\R^3$, and
$\mathcal G$ is open in $\mathrm{Sym}_3$.
Every joining path in a representing tree lies in $\mathcal G$.
Given $c_0,C_0$, the numbers $a,b$ can be chosen so that every
symmetric matrix with spectrum in $[c_0,C_0]$ belongs to
$\mathcal G$.
\end{lemma}

\begin{proof}
The product construction below is related to the diagonal
approximation in Marino and Spagnolo~\cite[Section~3]{MarinoSpagnolo};
we use its elementary coordinate-join arithmetic directly.
First we show that each scalar triple strictly between the endpoints
is interior. Given a nearby positive triple $(\alpha_1,\alpha_2,
\alpha_3)$, choose $M>\max_i\alpha_i$ and set
\[
 d_i=\sqrt{1-\alpha_i/M}.
\]
Let $t_i$ independently take the two values $1-d_i,1+d_i$ with
equal weights. Its arithmetic mean is one and its harmonic mean is
$1-d_i^2=\alpha_i/M$. Start with the eight scalar leaves
\begin{equation}
 \label{sc:product-leaves}
 M t_1t_2t_3\Id.
\end{equation}
Joining the $t_1$ pairs in direction 1 gives the diagonal triple
\[
 (\alpha_1t_2t_3,Mt_2t_3,Mt_2t_3).
\]
Joining the $t_2$ pairs in direction 2 gives
$(\alpha_1t_3,\alpha_2t_3,Mt_3)$, and the direction-3 joins give
$(\alpha_1,\alpha_2,\alpha_3)$. If the target tends to $(s,s,s)$,
choose $M$ tending to $s$ from above. Every leaf then tends to $s$,
so all leaves lie in $(a,b)$ for sufficiently close targets.

Interior propagates through any nontrivial join. With the other
child and a weight $0<\theta<1$ fixed, the derivative of the join
in its first child is diagonal, with entries $\theta$ in the
tangential positions and $\theta[J_i]_i^2/\alpha_i^2>0$ in the
normal position. It is invertible. The inverse function theorem,
applied successively up a finite tree, proves openness in diagonal
variables. Zero-weight joins can be removed. The diagonal class is
permutation invariant; continuity of ordered eigenvalues now proves
that its spectral extension is open, including at repeated
eigenvalues. Varying the weight of an existing join simply supplies
another finite tree with the same children, so its entire path,
including its endpoints, stays in $\mathcal G$.

For uniform coverage, take $M=2C_0$ in
Equation~\eqref{sc:product-leaves}. For targets in $[c_0,C_0]^3$,
all factors lie between
\[
 \ell=1-\sqrt{1-c_0/(2C_0)}>0
 \quad\text{and}\quad 2.
\]
All leaves therefore lie in $[M\ell^3,8M]$. For example,
$a=M\ell^3/2$ and $b=16M$ put them strictly inside the leaf
interval, uniformly for the entire required spectral range.
\end{proof}

Fix these $a,b$ for the rest of the proof. The distance of a joining
path from $\mathrm{Sym}_3\setminus\mathcal G$ can depend on the path,
but every fixed path is compact in the open set $\mathcal G$ and hence
has positive such distance. All corrections below are symmetric and
will respect that distance. Thus every corrected matrix remains in
the same $\mathcal G$, and the global bounds $a,b$ never deteriorate.

We now fix the equations and boundary data that every intermediate
field must preserve. Write $A^{\rm in},u_1^{\rm in},u_2^{\rm in}$
for the data of Theorem~\ref{thm:scalarization}, and put
$F_j^{\rm in}=A^{\rm in}\nabla u_j^{\rm in}$.

\begin{definition}[Subsolution]\label{sc:subsolution}
A subsolution consists of a measurable symmetric matrix $A$, two
potentials $u_1,u_2\in H^1(U)$, and the column matrices
\[
 E=(\nabla u_1,\nabla u_2),\qquad F=AE,
\]
such that $A\in\mathcal G$ almost everywhere and, for $j=1,2$,
\begin{align*}
 u_j-u_j^{\rm in}&\in H^1_0(U),& \Div F_j&=0,\\
 \int_U F_j\cdot\nabla\psi\,\dd x
 &=\int_U F_j^{\rm in}\cdot\nabla\psi\,\dd x
 &&\text{for every }\psi\in H^1(U).
\end{align*}
It must also have an open cover of a full-measure subset of $U$ on
which $A,u_1,u_2$ are smooth and have the local rank property in
Theorem~\ref{thm:scalarization}: on each member one fixed subset of
the potentials has independent gradients, and every remaining potential
is a constant affine combination of those potentials.
\end{definition}

The original fields are a subsolution by Lemma~\ref{sc:open-class}.
The construction will change the matrix and interior fields while keeping
the equations and the two entire boundary responses in this definition.

\subsection{A split in preliminary coordinates}

We first compute the states to be realized by a local perturbation.
A divergence-free flux wave with the potentials held fixed can mix the
tangential entries of a coordinate join. The normal entry requires its
harmonic mean instead. We obtain both at once by coupling the flux wave
to a normal change of variables; the following states describe that
coupling before localization.
Freeze a parent matrix $A_0$ and its two diagonal children $A^+$,
$A^-$ in a common orthonormal frame. Let the layering direction be
the unit vector $d$, and let the physical fractions be
$\theta_+,\theta_->0$, with sum one. A trivial join needs no
perturbation. Subscripts $d$ and $T$ denote the normal scalar entry
and tangential diagonal block. Set
\begin{equation}
 \label{sc:preliminary-states}
 q_s=\frac{A_d^s}{(A_0)_d},\qquad
 B_d^s=(A_0)_d,\qquad B_T^s=q_s A_T^s,
 \qquad \eta_s=\frac{\theta_s}{q_s},
 \quad s\in\{+,-\}.
\end{equation}
The normal harmonic-mean formula and the tangential arithmetic-mean
formula give
\begin{equation}
 \label{sc:means}
 \eta_++\eta_-=1,\qquad
 \eta_+q_++\eta_-q_-=1,\qquad
 \eta_+B^++\eta_-B^-=A_0.
\end{equation}
The fractions $\eta_s$, not $\theta_s$, are used before the
coordinate change.

Put $\Delta B=B^+-B^-$ and $\Delta q=q_+-q_-$. A scalar parameter
$z\in[-\eta_+,\eta_-]$ describes the segment
\begin{equation}
 \label{sc:synchronized-segment}
 B(z)=A_0+z\Delta B,\qquad q(z)=1+z\Delta q.
\end{equation}
Its endpoints are the two states in
Equation~\eqref{sc:preliminary-states}; its barycenter is $z=0$.
The normal stretch
$L(z)=\Id+(q(z)-1)d\otimes d$ pushes $B(z)$ forward to
\begin{equation}
 \label{sc:predicted-matrix}
 C(z)=\frac{L(z)B(z)L(z)^t}{\det L(z)}
      =\diag\bigl(q(z)(A_0)_d,\ B_T(z)/q(z)\bigr)
\end{equation}
in the chosen frame. This entire path is the original coordinate
joining path, with changed weights. Indeed if
$B=\sum_s\tau_s B^s$ and $q=\sum_s\tau_s q_s$ for
$\tau_++\tau_-=1$, the physical weights are
$w_s=\tau_s q_s/q$. Their normal harmonic mean is
\[
 \left(\sum_s\frac{w_s}{A_d^s}\right)^{-1}
 =q(A_0)_d,
\]
and their tangential arithmetic mean is $B_T/q$.
In particular, smoothing and cutting off the scalar parameter while
keeping it in its segment will not leave the admissible path.

\subsection{Exact localized flux waves}

Let $m\in\{1,2\}$ potentials have independent gradients on a
smooth patch. Temporarily discard the affinely dependent columns.
For a symmetric matrix to map their gradient matrix $E$ to a flux
matrix $F$, the pairing $E^tF$ must be symmetric. The flux wave must
preserve this condition as well as column divergence.
After shrinking the patch, choose smooth coordinates
\[
 z=\Phi(y),\qquad z_1=u_1(y),\ldots,z_m=u_m(y),
\]
where the independent potentials have been relabeled if needed.
If $J=D\Phi$, a flux column becomes the divergence density
$J F/\abs{\det J}$. In these coordinates, symmetry of $E^tF$
is exactly symmetry of the leading $m\times m$ block of this
transformed flux matrix. Thus the two constraints to preserve are
constant coefficient: column divergences vanish and that block is
symmetric.

Let $V_m$ be the vector space of $3\times m$ matrices with symmetric
leading block, equipped with the Frobenius inner product. Define
\[
 D(\xi):V_m\longrightarrow\R^m,\qquad D(\xi)G=\xi^tG.
\]
We identify the row vector on the right with a vector in $\R^m$.

The following is the surjective-symbol case of the polynomial
potential construction of Rai\textcommabelow{t}\u a~\cite[Theorem~1 and Section~2]{Raita}.
We give the explicit formula used for these coupled constraints.

\begin{lemma}[A polynomial potential]
\label{sc:polynomial-potential}
For every $\xi\ne0$, $D(\xi)$ is onto. There is a homogeneous
polynomial map $P(\xi):V_m\to V_m$ of degree $2m$ such that
\[
 D(\xi)P(\xi)=0,\qquad
 \operatorname{im}P(\xi)=\ker D(\xi)\quad(\xi\ne0).
\]
Consequently every mean-zero smooth periodic plane wave with
polarization in $\ker D(\xi)$ admits a compactly supported
localization satisfying both constraints exactly, with arbitrarily
small absolute uniform norm of the difference from its cut-off
leading wave.
\end{lemma}

\begin{proof}
For $m=1$, surjectivity is ordinary nonzero contraction. For $m=2$
write
\[
 G=\begin{pmatrix}a&b\\b&c\\d&e\end{pmatrix},\qquad
 D(\xi)G=(\xi_1a+\xi_2b+\xi_3d,
           \xi_1b+\xi_2c+\xi_3e).
\]
If $\lambda\in\R^2$ annihilates its range, then
\[
 \xi_1\lambda_1=\xi_2\lambda_2
 =\xi_1\lambda_2+\xi_2\lambda_1
 =\xi_3\lambda_1=\xi_3\lambda_2=0.
\]
When $\xi_3\ne0$ the assertion follows immediately; otherwise
either $\xi_1$ or $\xi_2$ is nonzero and the first three equalities
give the same conclusion. Thus $D(\xi)D(\xi)^t$ is positive
definite at every nonzero $\xi$, where the transpose uses the
specified inner products.

Set $Q(\xi)=D(\xi)D(\xi)^t$ and define
\begin{equation}
 \label{sc:polynomial-symbol}
 P(\xi)=\det Q(\xi)\Id
       -D(\xi)^t\operatorname{adj}(Q(\xi))D(\xi).
\end{equation}
The identity $Q\operatorname{adj}Q=(\det Q)\Id$ gives $DP=0$
as a polynomial identity. On $\ker D(\xi)$, $P(\xi)$ is
multiplication by the nonzero scalar $\det Q(\xi)$, proving the
image assertion. Both terms have degree $2m$.

For completeness, let $R\in\ker D(\xi)$ and choose
$R_0\in V_m$ with $P(\xi)R_0=R$. Let $h$ be a smooth
one-periodic function of mean zero and let $H$ be a one-periodic
$2m$-fold primitive, so $H^{(2m)}=h$. Such primitives are obtained
by successive integration and subtraction of the mean. Given a
smooth compactly supported cutoff $\chi$, apply the constant
coefficient differential operator $P(\partial)$ to
\begin{equation}
 \label{sc:localized-potential}
 k^{-2m}\chi(z)H(k\xi\cdot z)R_0.
\end{equation}
Its output lies in $V_m$ and has zero column divergence exactly,
by the polynomial identity. Terms with every derivative on $H$
give $\chi(z)h(k\xi\cdot z)R$. Each remaining term has at least
one derivative on $\chi$ and is $O(k^{-1})$ uniformly for fixed
$\chi,h$. This proves the localization claim.
\end{proof}

The flux wave will satisfy the differential constraints exactly but
will have a small constitutive error. The following algebraic fact
removes that error without losing symmetry.

\begin{lemma}[Symmetric algebraic correction]
\label{sc:symmetric-correction}
Let $E$ have independent columns and let $R$ satisfy that $E^tR$ is
symmetric. With $E^\dagger=(E^tE)^{-1}E^t$, the matrix
\begin{equation}
 \label{sc:correction-formula}
 H=R E^\dagger+(E^\dagger)^tR^t
       -(E^\dagger)^t(E^tR)E^\dagger
\end{equation}
is symmetric and satisfies $HE=R$. Its norm is bounded by a
locally uniform constant times $\norm{R}$ when $E$ ranges over a
bounded set with its smallest singular value bounded below.
\end{lemma}

\begin{proof}
Symmetry is immediate. Multiply by $E$, use $E^\dagger E=\Id$,
and cancel the last two terms using $R^tE=E^tR$. The norm bound
follows from the same formula.
\end{proof}

\subsection{Realizing one coordinate join}

We can now state the local operation used to exhaust a finite laminate
tree. Its two conclusions serve different purposes: proximity to the
children lowers the remaining tree depth, while weak proximity lets us
make that change inside any prescribed weak neighborhood.

\begin{lemma}[Local realization of a coordinate join]
\label{sc:local-join}
Let $(A,E,F)$, with potentials $u_1,u_2$, be a subsolution. Let $y_0$
lie in one of its smooth patches from Definition~\ref{sc:subsolution},
on which one fixed set of $m\in\{1,2\}$ potentials has independent
gradients. Suppose $A_0=A(y_0)$ is a coordinate join of
two children $A^+,A^-$ in one orthonormal frame, with positive weights
$\theta_+,\theta_-$ summing to one and with its joining path contained in
$\mathcal G$. Given open neighborhoods $\mathcal O_+,\mathcal O_-$ of
the children in $\mathrm{Sym}_3$ and $\delta>0$, there is an open
neighborhood $V$ of $y_0$, compactly contained in the patch, with the
following property.

For every open rectangular box $Q$ with $\overline Q\subset V$, one
can choose a cutoff and periodic waveform, then construct subsolutions
$(A_k',E_k',F_k')$, with potentials $u'_{1,k},u'_{2,k}$, for all
sufficiently large integers $k$ such that
$A_k'=A$ outside $Q$. The differences $u'_{j,k}-u_j$ and $F'_{j,k}-F_j$
vanish outside compact subsets of $Q$, and
\[
 \int_U(F'_{j,k}-F_j)\cdot\nabla\psi\,\dd x=0
 \qquad\bigl(\psi\in H^1(U),\ j=1,2\bigr).
\]
There are disjoint open subsets $Q_{k,+},Q_{k,-}\subset Q$ such that
\[
 A_k'\in\mathcal O_s\ \hbox{on }Q_{k,s}\quad(s\in\{+,-\}),\qquad
 \abs{Q\setminus(Q_{k,+}\cup Q_{k,-})}<\delta\abs Q.
\]
On the same fixed box, with the cutoff and waveform fixed,
\[
 (E_k',F_k')\rightharpoonup(E,F)
 \quad\text{in }L^2(U)\times L^2(U)\quad\text{as }k\to\infty.
\]
The neighborhood and frequency threshold may depend on all preceding
local data, including $\mathcal O_+,\mathcal O_-$ and $\delta$; the
threshold may also depend on the chosen box, cutoff, and waveform.
\end{lemma}

\begin{proof}
Choose an initial compactly contained neighborhood of $y_0$ on which
the potential chart $\Phi$ is defined and the selected gradients have
a positive lower singular-value bound. We will choose a smaller
neighborhood $V$ from the frozen-error estimates below; the box $Q$
will then have closure in $V$. Until the constitutive correction is
complete, the local matrices $E,F,E_0,R$ and their primed versions contain
only these $m$ independent columns. We restore the full pair by the fixed
affine relations at the end of the local calculation. Freeze
$A_0=A(y_0)$, $E_0=E(y_0)$, and $J_0=D\Phi(y_0)$.
For a split as above, the physical polarization
$\Delta B E_0$ has zero normal component since the normal row of
$\Delta B$ vanishes. Its gradient pairing is symmetric. Hence
\begin{equation}
 \label{sc:transformed-polarization}
 \xi=J_0^{-t}d,\qquad
 R=\frac{J_0\Delta B E_0}{\abs{\det J_0}}
 \quad\text{satisfy}\quad R\in V_m,\quad D(\xi)R=0.
\end{equation}
For the last assertion, contract the displayed matrix with $\xi$.
For the first, the leading rows of $J_0$ are the gradients of the
selected potentials, so their entries are precisely the symmetric
pairings with $\Delta B$.

The construction below applies on any box $Q$ with
$\overline Q\subset V$. After $V$ has been chosen from the estimates
below, fix such a box and a smooth cutoff $0\leq\chi\leq1$, supported
in $Q$ and equal to one except on a thin margin. Choose a smooth
periodic function $h$ by convolving the two-state
step function that equals $\eta_-$ on a fraction $\eta_+$ of
the period and $-\eta_+$ on the remaining fraction $\eta_-$.
Then
\begin{equation}
 \label{sc:wave-properties}
 \int_0^1h=0,\qquad -\eta_+\leq h\leq\eta_-,
\end{equation}
and $h$ equals the two pure values outside an arbitrarily small
part of its period. The margin and transition lengths will be chosen
to meet the prescribed volume error $\delta\abs Q$. Let $H_1$ be a
periodic primitive of $h$. In
Equation~\eqref{sc:localized-potential}, use the cutoff
$\chi\circ\Phi^{-1}$. Transform the resulting flux perturbation
back to the original coordinates and denote the perturbed flux by
$\widehat F_k$. It satisfies
\begin{equation}
 \label{sc:exact-preliminary-constraints}
 \Div\widehat F_k=0,\qquad
 E^t\widehat F_k=\widehat F_k^tE
\end{equation}
exactly, and equals the original flux outside a compact subset of
$Q$. The selected potentials have not yet changed.

Write $\varphi(y)=\xi\cdot\Phi(y)$; thus
$\nabla\varphi(y_0)=d$. At the same time define on $Q$, extending
by the identity outside $Q$,
\begin{equation}
 \label{sc:local-diffeomorphism}
 X_k(y)=y+\frac{\Delta q}{k}\,
             d\,\chi(y)H_1(k\varphi(y)).
\end{equation}
If $z_k(y)=\chi(y)h(k\varphi(y))$, differentiation gives
\begin{equation}
 \label{sc:jacobian-approximation}
 DX_k=\Id+\Delta q\,z_k\,d\otimes\nabla\varphi
          +O(k^{-1}).
\end{equation}
Throughout $Q$, this is uniformly close to $L(z_k)$ when $V$ is
sufficiently small and then $k$ is large. The scalar $z_k$ stays in the segment in
Equation~\eqref{sc:synchronized-segment}, because zero lies in
that segment and $0\leq\chi\leq1$. All $q(z_k)$ are bounded
away from zero. More explicitly, the rank-one determinant formula gives
\[
 \det DX_k=q(z_k)+\Delta q\,z_k(\partial_d\varphi-1)
       +\frac{\Delta q}{k}H_1(k\varphi)\partial_d\chi.
\]
First shrink $V$ so that $\partial_d\varphi$ is uniformly close to
one; after the box and profiles are fixed, take $k$ large. The displayed
determinant, which is the derivative of $X_k$ along each line parallel
to $d$ in that direction, is then strictly positive. The displacement is parallel
to $d$ and compactly supported. Each such line is consequently
mapped monotonically onto itself. This proves that $X_k$ is a
global smooth diffeomorphism, equal to the identity outside $Q$,
and mapping $Q$ onto itself. Its derivative and inverse derivative
have bounds independent of sufficiently large $k$.

Push the potentials and preliminary fluxes forward by this map:
\begin{equation}
 \label{sc:pushed-fields}
 u_{j,k}'=u_j\circ X_k^{-1},\qquad
 E_k'\circ X_k=DX_k^{-t}E,\qquad
 F_k'\circ X_k=\frac{DX_k\widehat F_k}{\det DX_k}.
\end{equation}
The determinant is positive here. These fields satisfy exact
column divergence and pairing symmetry, since
\[
 (E_k'^tF_k')\circ X_k
       =\frac{E^t\widehat F_k}{\det DX_k}.
\]
The selected gradients remain independent. The discarded potentials
are recomposed by the same map, so their original constant affine
relations with the selected potentials remain valid. Restore their
preliminary fluxes by the corresponding constant linear combinations
of the selected fluxes, and push them by the same formula. We continue
the constitutive estimates on the selected columns.

The errors in the frozen description are only errors in the
constitutive approximation, not in these equations. More precisely,
let $\omega(V)$ bound the oscillations from their values at $y_0$ of
the fixed smooth fields and chart derivatives on $V$, with
$\omega(V)\to0$ as $V$ shrinks to $y_0$. The frozen flux error can
be seen directly. Put $\Gamma(y)=\abs{\det J(y)}J(y)^{-1}$, so that
$\Gamma(y_0)R=\Delta B E_0$. The localized leading wave gives
\[
 \widehat F_k-B(z_k)E
 =(A-A_0)E+z_k\bigl[(\Gamma-\Gamma(y_0))R
                  +\Delta B(E_0-E)\bigr]+O(k^{-1}).
\]
Consequently, for the fixed plan,
\begin{align}
 \widehat F_k-B(z_k)E&=O(\omega(V))+O(k^{-1}),
       \label{sc:frozen-flux-error}\\
 DX_k-L(z_k)&=O(\omega(V))+O(k^{-1}).
       \label{sc:frozen-map-error}
\end{align}
All estimates here are absolute uniform bounds on the box. The
constants multiplying $\omega(V)$ depend only on the fixed plan and
bounds on the initial neighborhood; $z_k$ remains in its fixed segment.
The constants in $O(k^{-1})$ may also depend on the chosen box, cutoff,
and waveform. This is why those choices can be fixed before the frequency
is increased.

Define the constitutive defect
\[
 \mathcal R_k=F_k'-C(z_k\circ X_k^{-1})E_k'.
\]
Writing $K=DX_k$ and evaluating $B,L$ at $z_k(y)$, the pushforward
formulas give the exact identity
\[
 \mathcal R_k\circ X_k
 =\frac K{\det K}(\widehat F_k-BE)
 +\left(\frac{KBK^t}{\det K}-\frac{LBL^t}{\det L}\right)K^{-t}E.
\]
The map bounds and Equations~\eqref{sc:frozen-flux-error}--
\eqref{sc:frozen-map-error} therefore give
\begin{equation}
 \label{sc:constitutive-error}
 \mathcal R_k=O(\omega(V))+O(k^{-1}).
\end{equation}

Apply Lemma~\ref{sc:symmetric-correction} to $\mathcal R_k$, using
the independent columns. Its symmetry condition is exact: both
$E_k'^tF_k'$ and $E_k'^tC(z_k\circ X_k^{-1})E_k'$ are symmetric.
The selected gradients retain a positive lower singular-value bound,
because $E_k'\circ X_k=DX_k^{-t}E$ and the map derivatives are bounded.
Thus the correction $H_k$ has a uniform bound by a local constant
times $\abs{\mathcal R_k}$, and
\[
 A_k'=C(z_k\circ X_k^{-1})+H_k
\]
maps the selected new gradients to their new fluxes. Restore every
discarded flux by the same constant linear combination as its gradient.
The common recomposition then gives the full two-column relation
$F_k'=A_k'E_k'$ and preserves the original fields outside the box.
From this point onward, $E,F,E_k',F_k'$ and $\widehat F_k$ again denote
the full two-column matrices; the bound for $H_k$ is the one just obtained
from the selected columns.

Choose $\tau>0$ so that the symmetric $\tau$-neighborhood of the compact
path $C$ lies in $\mathcal G$ and the $\tau$-balls about $A^s$ lie in
$\mathcal O_s$ for $s\in\{+,-\}$. This uses distance inside $\mathrm{Sym}_3$.
The coefficient of $\omega(V)$ in the correction bound is fixed on the
initial neighborhood, independently of the later box and profiles.
Choose $V$ small enough that this contribution is less than $\tau/2$
and that the preceding map determinant is positive after a sufficiently
small frequency error. For each box with closure in this $V$, fix its
cutoff and waveform and take $k$ large enough that the remaining
contribution is less than $\tau/2$. Then $A_k'\in\mathcal G$ throughout
the box. The correction needs no bounds on derivatives of the defect.
It is smooth inside the box; outside $Q$ keep the original matrix.
A jump of the matrix at the box boundary is harmless: the potentials
and fluxes agree with the old ones near that boundary, and the matrix
is only required to be smooth on an open cover up to null sets.

Where $\chi=1$ and $h$ is at the pure value for child $s$, the matrix
$C$ is exactly $A^s$, so the correction bound puts $A_k'$ in
$\mathcal O_s$. Take $Q_{k,s}$ to be the images under $X_k$ of the
open pure regions, omitting the endpoints of the waveform plateaus.
The omitted plateau endpoints form a null set. The cutoff margin can
have arbitrarily small volume, and the transition intervals can have
arbitrarily small total length. The phase is nondegenerate on the box;
changing to the chart and integrating in its linear phase direction
shows that the preimages of the transition intervals have the
corresponding small limiting volume. For the image measure, the
determinant formula above gives
\[
 0<\det DX_k\le
 \sup_{z\in[-\eta_+,\eta_-]}q(z)
 +\abs{\Delta q}\sup_{z\in[-\eta_+,\eta_-]}\abs z
       \sup_V\abs{\partial_d\varphi-1}+1=:J_*,
\]
after the profiles are fixed and $k$ is large enough to make the $k^{-1}$
term in the determinant formula at most one in absolute value. Thus $J_*$ is
fixed by the preceding local data and $V$, independently of the cutoff margin and
transition lengths, and bounds the volume increase under $X_k$.
Choose the margin and transition lengths so that, for all sufficiently large $k$, the
complement of $Q_{k,+}\cup Q_{k,-}$ has volume less than
$\delta\abs Q$. The new smooth box pieces and the old pieces away
from its faces retain the full-measure local rank cover.

\paragraph{Boundary data and weak proximity.}

Every perturbation just constructed preserves the outer boundary
traces and fluxes exactly. Before the pushforward the flux change
is divergence-free and compactly supported in the interior. After
the pushforward the same is true of its difference from the old
flux, since $X_k$ is the identity outside $Q$. Thus for every
$\psi\in H^1(U)$ the flux difference has zero pairing with
$\nabla\psi$. To see this without any boundary smoothness of the
fields, insert a smooth cutoff equal to one on the support of the
flux difference, and use density of smooth tests in $H^1_0(U)$ together
with its distributional zero divergence. The potentials also agree
outside a compact subset of $Q$, so their trace difference vanishes.
There is no new global boundary compatibility
condition: the perturbation starts from an existing solution pair
and satisfies the exact local equations with compact support.
Together with the constitutive relation and the local rank cover, these
facts prove that the corrected fields are subsolutions.

It is equally important that these perturbations can be made
arbitrarily weakly close to the old fields, even though the local
child states stay separated. Keep the chosen $V$, box, waveform,
and cutoff fixed. The preliminary flux
perturbation is a smooth matrix amplitude times
$h(k\varphi(y))$, plus a uniformly $O(k^{-1})$ error.
Changing to the chart and using the zero mean of $h$ shows
\begin{equation}
 \label{sc:preliminary-weak-limit}
 \widehat F_k\rightharpoonup F\quad\text{in }L^2.
\end{equation}
For example, for a smooth compactly supported test amplitude one
can integrate a bounded periodic primitive in the phase direction,
obtaining $O(k^{-1})$; density and the uniform $L^2$ bound give
the stated weak convergence.

The product of the oscillating Jacobian and flux in
Equation~\eqref{sc:pushed-fields} cannot be treated by multiplying
weak limits. Instead write $v_k=X_k-\Id$. For a smooth vector
test function $\psi$, change variables to obtain
\[
 \int_U F_k'(x)\cdot\psi(x)\,\dd x
 =\int_U DX_k(y)\widehat F_k(y)\cdot\psi(X_k(y))\,\dd y
\]
for each flux column. Replacing $\psi\circ X_k$ by $\psi$
costs $o(1)$ because $v_k\to0$ uniformly and the other factors
are uniformly bounded in $L^2$. The remaining additional term
vanishes by the exact divergence equation:
\begin{equation}
 \label{sc:piola-weak-cancellation}
 \int_U\partial_j(v_k)_i(\widehat F_k)_j\psi_i\,\dd y
 =-\int_U(v_k)_i(\widehat F_k)_j\partial_j\psi_i\,\dd y
 \longrightarrow0.
\end{equation}
Summation over repeated spatial indices is understood only in this
display. Equations~\eqref{sc:preliminary-weak-limit} and
\eqref{sc:piola-weak-cancellation} prove $F_k'\rightharpoonup F$.
The potentials converge uniformly on the modified compact region,
since $X_k^{-1}\to\Id$ uniformly and the old potentials are
smooth there. Their gradients are uniformly bounded in $L^2$,
either directly on the box or by the energy estimate below.
Consequently $E_k'\rightharpoonup E$ as well. The same weak limits hold
for the restored columns, since their local changes are the same fixed
linear combinations of the selected-column changes.

The choices now have the order asserted in the lemma. First the
neighborhood $V$ controls the fixed coefficient error $O(\omega(V))$
for every box whose closure lies in $V$. After fixing one such box,
choose its cutoff and waveform to meet the volume budget. All
sufficiently large frequencies then give the pure-state membership,
and that same family converges weakly to the old pair as the frequency
tends to infinity. The freezing error is not an error in the
differential constraints and does not obstruct this limit.
\end{proof}

\subsection{Finite-depth exhaustion}

We use Lemma~\ref{sc:local-join} to replace a finite laminate tree by
its children one depth at a time. Every step remains a subsolution in
the explicit sense of Definition~\ref{sc:subsolution}.

\begin{proposition}[Weak approximation by scalar-near matrices]
\label{sc:weak-approximation}
Given a subsolution, a weak $L^2$ neighborhood of its pair $(E,F)$,
and $\eps>0$, there is a subsolution in that neighborhood whose
matrix has distance less than $\eps$ from the scalar matrices
outside a set of volume less than $\eps$.
\end{proposition}

\begin{proof}
For clarity we make the finite-depth argument explicit. In diagonal
variables let $T_0$ be the triples in $\mathcal G^{\mathrm{diag}}$
whose distance from scalar triples is strictly less than $\eps$.
Inductively, $T_k$ contains $T_{k-1}$ and the outputs of joins of
two $T_{k-1}$ children whose entire joining path lies in
$\mathcal G^{\mathrm{diag}}$. Equivalently these are targets
with a representing tree of depth at most $k$, ending in $T_0$,
with the indicated path condition. Take their spectral extensions
when discussing matrices.

Each $T_k$ is open and permutation invariant. This is immediate
for $T_0$. For a nontrivial join with children in $T_{k-1}$, its
invertible derivative in one child, computed in
Lemma~\ref{sc:open-class}, absorbs any sufficiently small change
of output into a change of that child. The child remains in the
open set $T_{k-1}$, and the whole path remains in the open class
by compactness. This proves the induction; permutation invariance
and spectral openness follow as before. Every member of
$\mathcal G$ lies in some $T_k$, since its defining tree ends
at scalar leaves.

Intersect the smooth rank patches with the set where $A\in\mathcal G$.
This intersection is open in each patch and still has full measure.
Its subsets where $A\in T_k$ are open and increase to the whole
intersection. By finite measure, choose a finite $K$ so that the portion
not covered by these rank-regular $T_K$ patches has volume less than
$\eps/3$. If $K=0$, the original subsolution already suffices. Otherwise
fix $\delta>0$ with $\delta\abs U<\eps/(3K)$. At each of the $K$
reductions, take a compact subset of the current good region losing
volume less than $\eps/(3K)$, and use this fixed relative tolerance
$\delta$ in the local lemma before choosing its neighborhoods.
Consider one reduction from $T_k$ to $T_{k-1}$, with $1\leq k\leq K$.
Choose every admissible neighborhood below inside its current open
$T_k$ rank patch.

At a rank-positive point, fix its smooth potential chart. If its
matrix already belongs to $T_{k-1}$, choose an admissible neighborhood
on which $A\in T_{k-1}$; every contained box then needs no change.
Otherwise choose its frozen two-child
plan, with children in $T_{k-1}$. Apply Lemma~\ref{sc:local-join}
with both child neighborhoods equal to the spectral extension of
$T_{k-1}$ and with the already fixed relative tolerance $\delta$.
It supplies an admissible neighborhood for every contained grid box;
the large-frequency family then permits any prescribed weak error.
The local independent gradients
have a positive smallest singular value on the initial compact
neighborhood used by that lemma. At a rank-zero point, choose an
admissible neighborhood on which the potentials are constant and their
fluxes vanish. On any box with closure inside it, replace the matrix by
any scalar strictly in $(a,b)$, without changing the fields.

These choices give admissible neighborhoods of every point in the
chosen compact subset. A finite subcover of that compact set and a fine
enough grid give finitely many disjoint box interiors, with
closures inside admissible neighborhoods, covering the compact
set except for grid faces. The frozen point need not lie in the
grid box: Lemma~\ref{sc:local-join} applies to every box whose closure
lies in its admissible neighborhood. One can equivalently refine a grid
successively. Each box is treated using its own chart, fixed
frame, rank bound, and frozen plan; no continuous eigenframe on
the whole domain is required.

For this reduction, construct every boxwise output from the same incoming
subsolution, with the frozen data just selected. In each rank-positive
box requiring a join, use the local lemma with the fixed relative tolerance
$\delta$. Paste the individual potential and flux differences and the
matrix choices over the disjoint box interiors. The potential and flux
differences have compact support in those boxes, so the pasted gradients and fluxes have each box's
constitutive relation, and the exact flux pairings add. Since the box
interiors are disjoint, the sum of the exceptional volumes is less than
$\delta\sum_Q\abs Q\leq\delta\abs U<\eps/(3K)$. The unchanged
$T_{k-1}$ boxes and the rank-zero scalar boxes create no exceptional
volume. Together with them, the open pure regions give smooth patches
with the local rank property on which the new matrix belongs to $T_{k-1}$.
The output is again a subsolution. In joined boxes this is part of the
lemma; in the other boxes the fields are unchanged, and either the
matrix is unchanged or it maps the zero gradients to the zero fluxes.
The scalar replacement lies in $\mathcal G$ and retains the local rank
property. Matrix jumps at the finitely many box faces add only null seams,
and the original cover remains where nothing changed.

Perform the next step only inside these new good patches and
continue down to $T_0$. Untreated exceptional portions can be
left untouched. Every map is supported in a selected box and
maps that box onto itself. Thus later operations can be confined
to the current good set. The initial loss is less than $\eps/3$;
the $K$ compact-subset losses and the $K$ local exceptional losses
are each less than $\eps/(3K)$ per stage. Their total is less than
$\eps$. Fixed global
rank bounds or fixed tolerances over infinitely many plans are
not needed.

Finally, because the box outputs are pasted from one incoming pair, the
finite sum of their field differences is arbitrarily weakly small by taking
their frequencies sufficiently large. Use a metric for the weak topology on the
bounded set of all subsolution pairs, whose boundedness is shown
below, and allot a sufficiently small part of the prescribed
neighborhood to each of the finitely many stages. The final pair
lies in the prescribed neighborhood, and its matrix is in $T_0$
off a set of measure less than $\eps$, as asserted.
\end{proof}

\subsection{The strong limit selected by Baire}

We now close the argument without losing either Cauchy condition.
All subsolution pairs are uniformly bounded in
\[
 \mathcal H=L^2(U;\R^{3\times2})\times L^2(U;\R^{3\times2}).
\]
Indeed, for each potential $\widetilde u_j$ of a subsolution,
energy minimality with its fixed boundary trace gives
\begin{equation}
 \label{sc:uniform-energy}
 a\int_U\abs{\nabla\widetilde u_j}^2
 \leq\int_U \widetilde A\nabla\widetilde u_j\cdot\nabla\widetilde u_j
 \leq b\int_U\abs{\nabla u_j^{\rm in}}^2.
\end{equation}
Here $\widetilde A$ is that subsolution's matrix.
Also $\abs{F}\leq b\abs{E}$. Since
$\widetilde u_j-u_j^{\rm in}\in H^1_0(U)$, the zero-trace Poincar\'e
inequality bounds the potentials themselves as well.

Let $Z$ be the weak closure in $\mathcal H$ of all subsolution
pairs. It is nonempty, weakly compact, and metrizable in its weak
topology. One may use a metric obtained from a countable
orthonormal basis of the separable Hilbert space $\mathcal H$.
Every point of $Z$ still consists of gradients of potentials with
the prescribed traces and of divergence-free fluxes with the
prescribed normal-flux functionals. To verify the gradient
assertion, take a weakly convergent sequence of subsolutions;
Equation~\eqref{sc:uniform-energy} and Poincar\'e give weakly
convergent potentials in the fixed affine space
$u_j^{\rm in}+H^1_0(U)$. Their weak gradients are the limiting columns.
For the flux assertion, all functionals
$F_j\mapsto\int_U F_j\cdot\nabla\psi$ are weakly continuous,
for every fixed $\psi\in H^1(U)$.

\begin{lemma}[Weak-to-strong continuity points]
\label{sc:baire}
The identity from $Z$ with its weak topology to $\mathcal H$ with
its norm topology has a dense set of continuity points.
\end{lemma}

\begin{proof}
Let $P_N$ be the finite-rank orthogonal projections onto the first
$N$ vectors of a fixed orthonormal basis. Their restrictions to $Z$
are weak-to-norm continuous and $P_Nz\to z$ in norm for every
$z$. Here is the relevant Baire argument explicitly. On any
nonempty closed subset $K$ of $Z$ and for $\delta>0$, the closed
sets
\[
 K_N=\{z\in K:\norm{P_pz-P_qz}\leq\delta
                    \text{ for all }p,q\geq N\}
\]
cover $K$. Compact metric spaces are Baire, so some $K_N$ has
nonempty relative interior. On that patch, letting $p$ tend to
infinity gives $\norm{z-P_Nz}\leq\delta$. Shrinking the patch
using continuity of $P_N$, its strong diameter is at most
$3\delta$.

Apply this inside the closure of an open set whose closure lies
in any prescribed nonempty open subset of $Z$. The relative patch
meets the original open set, because that set is dense in its
closure. Their intersection contains a nonempty open subset of $Z$
of strong diameter at most $3\delta$. Consequently, for each
$r>0$, the union of all open subsets of $Z$ of strong diameter
less than $r$ is open and dense. The intersection of these sets
for $r=1,1/2,1/3,\ldots$ is dense by Baire, and every point of
the intersection is a continuity point of the identity.
\end{proof}

\begin{proof}[Proof of Theorem~\ref{thm:scalarization}]
Choose a continuity point $(E,F)\in Z$ from
Lemma~\ref{sc:baire}. By the definition of $Z$, choose subsolution
pairs converging weakly to it, and apply
Proposition~\ref{sc:weak-approximation} to each, with weak errors
tending to zero and with $\eps_j=2^{-j}$. Denote the resulting
subsolutions by $(A_j,E_j,F_j)$. Then
\begin{equation}
 \label{sc:strong-continuity-limit}
 (E_j,F_j)\longrightarrow(E,F)\quad\text{strongly in }\mathcal H,
\end{equation}
and $\dist(A_j,\R\Id)<2^{-j}$ except on a set of measure less
than $2^{-j}$. Put $\lambda_j=\tr(A_j)/3$. Since $A_j\in
\mathcal G$, one has $a<\lambda_j<b$, and the Frobenius-distance
projection onto scalar matrices is exactly $\lambda_j\Id$.

The exceptional measures are summable. After discarding their
limsup, which has measure zero, we have
$A_j-\lambda_j\Id\to0$ pointwise. Passing to a subsequence in
Equation~\eqref{sc:strong-continuity-limit} also gives pointwise
convergence of $E_j,F_j$ almost everywhere. At such a point,
\[
 F_j-\lambda_jE_j=(A_j-\lambda_j\Id)E_j\longrightarrow0.
\]
If $E\ne0$, it follows that
\begin{equation}
 \label{sc:measurable-scalar}
 \lambda_j\longrightarrow\gamma=
       \frac{E:F}{\abs{E}^2}\in[a,b],\qquad F=\gamma E,
\end{equation}
where $E:F$ is the Frobenius pairing. If $E=0$, the bound
$\abs{F_j}\leq b\abs{E_j}$ forces $F=0$; define $\gamma=a$
there. This defines one measurable scalar, common to both columns,
with $a\leq\gamma\leq b$ almost everywhere. No pointwise
subsequence of the bounded coefficient sequence itself was
assumed.

The strong convergence is essential for excluding concentration
on small exceptional sets. Equivalently, it makes $\abs{E_j}^2$
uniformly integrable, so the coefficient defect times $E_j$
also tends to zero in $L^2$. A mere uniform energy bound would
not suffice for that conclusion.

Finally the gradient and flux characterization of $Z$ supplies
$v_1,v_2\in H^1(U)$ with the original traces, gradients $E$,
and the original normal-flux functionals. Since $F=\gamma E$
and its columns have zero divergence, these are precisely
Equations~\eqref{sc:conclusion-equations} and
\eqref{sc:conclusion-flux}.
\end{proof}
